%% sections/02-abstract.tex
%% Abstract & Purpose
%% IEC/IEEE 82079-1:2019 — Conceptual Information
%% FAIR: Findable (keywords, structured summary)
%% ─────────────────────────────────────────────────

\section*{Abstract}
\addcontentsline{toc}{section}{Abstract}
\label{sec:abstract}

%%
%% WRITING GUIDE — ABSTRACT (3 structured fields):
%%
%%   WHAT THIS NOTE COVERS: One sentence. Name the task and the system.
%%   WHAT YOU WILL ACHIEVE: The concrete outcome after following this note.
%%   SCOPE: What is included and explicitly what is excluded.
%%
%% Keep the total abstract under 200 words.
%%

\noindent
\begin{tcolorbox}[
  enhanced, breakable,
  colback=PBALightGray, colframe=PBABlue,
  leftrule=4pt, rightrule=0.5pt, toprule=0.5pt, bottomrule=0.5pt,
  arc=2pt, left=10pt, right=10pt, top=8pt, bottom=8pt,
]

\textbf{What this note covers:}\enspace
This application note describes the complete procedure for configuring and
commissioning force control mode on a PBA Systems GS-Series gantry system
using an ACS SPiiPlus motion controller.

\medskip
\textbf{What you will achieve:}\enspace
After following this note you will have a stable, calibrated force control
loop that can apply a defined normal force in the Z-axis during bonding or
pick-and-place tasks, with live force feedback visible on the ACS SPiiPlus
MMI and logged via Python.

\medskip
\textbf{Scope:}\enspace
This note covers sensor wiring, ACS controller parameter setup, a Python
logging script, verification of the closed-loop response, and a
troubleshooting table.
It does \textit{not} cover: mechanical sensor installation, network
configuration of the ACS controller, or process-specific force profiles.
Refer to the relevant mechanical drawing and \mbox{PBA-AN-002-2025} for
those topics.

\end{tcolorbox}

\clearpage
